\documentclass[12pt,a4paper]{article}

% --- CÀI ĐẶT FONT CHỮ VÀ NGÔN NGỮ ---
\usepackage{fontspec}
\setmainfont{Times New Roman}
\usepackage{amsmath}

% --- CĂN LỀ TRANG ---
\usepackage[top=2cm, bottom=2cm, left=2.5cm, right=2cm]{geometry}

% --- LÀM ĐẸP ĐOẠN VĂN (PARAGRAPH) ---
\usepackage{parskip} % Tự động tạo khoảng cách giữa các đoạn, bỏ lùi đầu dòng
\usepackage{microtype} % Căn chữ (justify) đều và đẹp hơn, tránh bị giãn khoảng trắng

% --- LÀM ĐẸP TIÊU ĐỀ (HEADINGS) ---
\usepackage{titlesec}
\usepackage{xcolor}
\titleformat{\section}{\Large\bfseries\color{blue!50!black}}{\thesection.}{1em}{}
\titleformat{\subsection}{\large\bfseries\color{blue!60!black}}{\thesubsection.}{1em}{}
\titleformat{\subsubsection}{\normalsize\bfseries\color{blue!70!black}}{\thesubsubsection.}{1em}{}

% --- LÀM ĐẸP BẢNG & HÌNH ẢNH ---
\usepackage{booktabs}
\usepackage{longtable}
\usepackage{array}
\usepackage{graphicx}
\usepackage{float}
\usepackage{caption}
\captionsetup{labelfont=bf, textfont=it, margin=1cm}

% --- HYPERLINK ---
\usepackage[colorlinks=true, linkcolor=blue, urlcolor=blue, citecolor=blue]{hyperref}

% Định nghĩa lại pandocbounded nếu pandoc 3.x sử dụng nó
\providecommand{\pandocbounded}[1]{#1}

% Định nghĩa \tightlist để tránh lỗi biên dịch list của Pandoc
\providecommand{\tightlist}{\setlength{\itemsep}{0pt}\setlength{\parskip}{0pt}}

% Định nghĩa counter giả 'none' để sửa lỗi tương thích longtable/caption của Pandoc
\newcounter{none}

\begin{document}


\section{BÁO CÁO TỔNG THỂ CÁC MÔ HÌNH NHẬN DIỆN TÉ NGÃ (v30 -
v31)}\label{buxe1o-cuxe1o-tux1ed5ng-thux1ec3-cuxe1c-muxf4-huxecnh-nhux1eadn-diux1ec7n-tuxe9-nguxe3-v30---v31}

\textbf{Mục tiêu:} Tổng hợp, phân tích và so sánh các kiến trúc mạng
Neural Network từ phiên bản v30 đến v31 được tối ưu hóa (TinyML) để chạy
trên vi điều khiển ESP32-S3 sử dụng bộ gia tốc ESP-NN.

\begin{center}\rule{0.5\linewidth}{0.5pt}\end{center}

\subsection{1. DỮ LIỆU VÀ TIỀN XỬ LÝ (DATA
PIPELINE)}\label{dux1eef-liux1ec7u-vuxe0-tiux1ec1n-xux1eed-luxfd-data-pipeline}

Tất cả các phiên bản đều sử dụng chung một cơ sở dữ liệu SisFall với cấu
hình Pipeline được thiết kế vô cùng tinh xảo: - \textbf{Phân loại 5 nhãn
(Classes):} \texttt{Walk}, \texttt{Run}, \texttt{Idle}, \texttt{Trans},
\texttt{Fall} (Té ngã là index 4).

\subsubsection{1.1. Tiền xử lý \& Giảm tần số lấy mẫu
(Downsampling)}\label{tiux1ec1n-xux1eed-luxfd-giux1ea3m-tux1ea7n-sux1ed1-lux1ea5y-mux1eabu-downsampling}

\begin{itemize}
\tightlist
\item
  Tần số lấy mẫu gốc của cảm biến SisFall là 200Hz.
\item
  Thuật toán tiến hành \textbf{Downsampling bằng cách lấy mẫu xen kẽ
  (\texttt{df.iloc{[}::2,\ :{]}})} để giảm tần số xuống \textbf{100Hz}.
  Việc này giúp giảm một nửa khối lượng tính toán mà vẫn giữ nguyên được
  các đặc trưng tần số quan trọng của hành vi con người.
\item
  \textbf{Chuẩn hoá (Normalization):}

  \begin{itemize}
  \tightlist
  \item
    Gia tốc chia 8 (kẹp \texttt{{[}-8,\ 8{]}\ g}) để đưa về dải
    \texttt{{[}-1,\ 1{]}}.
  \item
    Gyro chia 500 (kẹp \texttt{{[}-500,\ 500{]}\ dps}) đưa về dải
    \texttt{{[}-1,\ 1{]}}. Giới hạn 500 dps được chọn dựa trên phân tích
    bách phân vị (p99.9) của tập dữ liệu thay vì lấy mốc 2000 dps gây
    lãng phí độ phân giải khi lượng tử hoá INT8.
  \item
    \textbf{Lưu ý Lịch sử (Bug ở v25):} Ở phiên bản \texttt{v25}, dữ
    liệu Gyro bị nhầm đơn vị là \texttt{rad/s} thay vì \texttt{dps}. Do
    1 rad/s ≈ 57.3 dps, toàn bộ biên độ Gyro đã bị thu nhỏ đi 57 lần.
    Việc \texttt{v25} tiếp tục lấy giá trị này chia cho 2000 khiến dữ
    liệu Gyro khi lượng tử hoá về \texttt{INT8} (dải -128 đến 127) trở
    nên cực kỳ bé, gần như quy về 0. Hệ quả: Mô hình \texttt{v25} ``mù''
    Gyro, chỉ dựa vào gia tốc (Accel) để đánh giá. Việc mất đi 3 trục
    Gyro vô tình loại bỏ một lượng lớn tín hiệu nhiễu phức tạp, giúp
    mạng Neural \textbf{học và hội tụ dễ dàng hơn rất nhiều} trong quá
    trình train. Tuy nhiên, cái giá phải trả là mô hình mất đi hoàn toàn
    khả năng cảm nhận không gian xoay, làm giảm hẳn sự nhạy bén thực tế
    (Recall) đối với những hành động phức tạp.
  \end{itemize}
\end{itemize}

\subsubsection{1.2. Kỹ thuật Cắt cửa sổ (Windowing) \& Gán nhãn nâng
cao}\label{kux1ef9-thuux1eadt-cux1eaft-cux1eeda-sux1ed5-windowing-guxe1n-nhuxe3n-nuxe2ng-cao}

Thay vì cắt Sliding Window mù quáng, hệ thống áp dụng kỹ thuật
\textbf{Cắt theo sự kiện (Event-based Windowing)} cực kỳ thông minh: -
\textbf{Kích thước cửa sổ:} 200 mẫu (tương đương 2 giây ở tần số 100Hz).
- \textbf{Với nhóm hành động chuyển tiếp (Trans):} Hệ thống dùng hàm
\texttt{rolling-RMS} của Gyro với ngưỡng kích hoạt
\texttt{\textgreater{}\ 20\ dps}. Các vùng liên tục được gom thành 1 sự
kiện và cắt duy nhất 1 cửa sổ 2 giây căn đúng vào đỉnh của sự kiện đó.
Kỹ thuật này giúp mô hình chống nhiễu (robust) cực tốt kể cả khi hành
động có nhiều đỉnh phụ (sub-peaks). - \textbf{Phân tách nhãn thông
minh:} Những đoạn dữ liệu nằm xa sự kiện Gyro được gán nhãn
\texttt{Idle}. - \textbf{Xử lý nhiễu đặc biệt (D18 \& D19):} Các mẫu D18
(vấp khi đi) và D19 (nhảy nhót) chứa rất nhiều nhiễu. Thuật toán chủ
động chỉ trích xuất đúng 1 cửa sổ \texttt{Trans} tại đỉnh gia tốc (SVM)
cao nhất, toàn bộ phần đi bộ xung quanh bị vứt bỏ để tránh làm ``bẩn''
nhãn \texttt{Walk} hay \texttt{Idle}. - \textbf{Sliding Window truyền
thống:} Chỉ áp dụng cho các hành động mang tính chu kỳ liên tục như
\texttt{Walk} (D01, D02, D05) và \texttt{Run} (D03, D04).

\subsubsection{1.3. Data Augmentation}\label{data-augmentation}

\begin{itemize}
\tightlist
\item
  Dữ liệu thuộc lớp \texttt{Trans} được làm phong phú (Augment) trong
  quá trình huấn luyện bằng cách \textbf{nhân biên độ tỷ lệ
  \texttt{x0.9} và \texttt{x1.1}} thay vì dịch chuyển cửa sổ (sliding
  shift). Điều này giúp mô hình học được sự biến thiên lực của các đối
  tượng (người lớn tuổi/người trẻ) có vóc dáng khác nhau.
\end{itemize}

\subsubsection{1.4. Cấu trúc Tensor Đầu
Vào}\label{cux1ea5u-truxfac-tensor-ux111ux1ea7u-vuxe0o}

\begin{itemize}
\tightlist
\item
  \textbf{Từ v30 đến v30\_tcn\_optimize\_v2 (6 trục):}
  \texttt{{[}1,\ 200,\ 6{]}} (\(a_x\), \(a_y\), \(a_z\), \(g_x\),
  \(g_y\), \(g_z\)).
\item
  \textbf{Đặc biệt với v31 (4 trục):} Rút gọn đầu vào thành \(a_x\),
  \(a_y\), \(a_z\) và \textbf{Gyro RMS}. Gyro RMS được tính bằng
  \(\text{Gyro}_{\text{RMS}} = \sqrt{g_x^2 + g_y^2 + g_z^2}\). Tensor
  đầu vào là \texttt{{[}1,\ 200,\ 4{]}}. Trên Firmware ESP32, vi điều
  khiển sẽ tự tính toán \(\text{Gyro}_{\text{RMS}}\) realtime từ 6 trục
  UART.
\end{itemize}

\subsubsection{1.5. Biêu đồ Phân bố Dữ liệu (Data Distribution \&
Thresholds)}\label{biuxeau-ux111ux1ed3-phuxe2n-bux1ed1-dux1eef-liux1ec7u-data-distribution-thresholds}

Dưới đây là các biểu đồ Histogram (Log Base 10 Scale) minh hoạ trực quan
toàn bộ phân bố của các trục tín hiệu gốc kèm vạch đỏ đánh dấu ngưỡng
cắt (Threshold) tối ưu và vạch trọng lực Trái Đất (Gravity):

\textbf{1. Phân bố Gia tốc (Accelerometer Distribution 2x2)}
\emph{(Layout 2x2 hiển thị phân bố các trục X, Y, Z và SVM tổng hợp.
Vạch đỏ thể hiện ngưỡng cắt tối ưu 8.0g, vạch đen đứt thể hiện trọng lực
Trái Đất 1.0g)}
\begin{figure}[H]
    \centering
    \pandocbounded{\includegraphics[width=0.85\textwidth, keepaspectratio,alt={Biểu đồ phân bố các trục gia tốc layout 2x2}]{./accel_distribution_2x2.png}}
\end{figure}

\newpage

\textbf{2. Phân bố Vận tốc góc (Gyroscope Distribution 2x2)}
\emph{(Layout 2x2 hiển thị phân bố các trục X, Y, Z và RMS tổng hợp.
Vạch đỏ thể hiện ngưỡng cắt tối ưu 500 dps)}
\begin{figure}[H]
    \centering
    \pandocbounded{\includegraphics[width=0.85\textwidth, keepaspectratio,alt={Biểu đồ phân bố các trục vận tốc góc layout 2x2}]{./gyro_distribution_2x2.png}}
\end{figure}

\textbf{Bảng Phân Tích Bách Phân Vị (Percentiles - Tính theo giá trị
tuyệt đối):}

Bảng dưới đây thống kê các mốc phân bố tín hiệu nhằm làm cơ sở thiết lập
thang đo phần cứng (8g và 500 dps) để lượng tử hóa (Quantization INT8):

{\def\LTcaptype{none} % do not increment counter
\begingroup
\small
\setlength{\tabcolsep}{4pt}
\renewcommand{\arraystretch}{1.3}
\begin{longtable}{@{} l c c c @{}}
\toprule\noalign{}
Trục Tín Hiệu & P50 (Trung vị) & P95 & P99.9 \\
\midrule\noalign{}
\endhead
\bottomrule\noalign{}
\endlastfoot
Accel X (g) & 0.09 & 0.98 & 2.16 \\
Accel Y (g) & 0.90 & 1.44 & 4.49 \\
Accel Z (g) & 0.27 & 0.97 & 2.62 \\
Accel SVM (g) & 1.00 & 1.60 & 5.22 \\
Gyro X (dps) & 4.39 & 70.13 & 404.49 \\
Gyro Y (dps) & 4.58 & 65.12 & 240.12 \\
Gyro Z (dps) & 1.77 & 49.93 & 227.13 \\
Gyro RMS (dps) & 11.39 & 111.17 & 463.57 \\
\end{longtable}
\endgroup
}

\emph{Nhận xét:} - \textbf{Gia tốc (Accel):} Ở bách phân vị 99.9\%, mức
gia tốc lớn nhất (của trục Y và SVM) chỉ xoay quanh mốc 4.5g đến 5.2g.
Việc chọn thang đo phần cứng \texttt{±8g} bao phủ hoàn toàn mọi dao động
khốc liệt nhất của các cú ngã mà không bị cắt xén dữ liệu (clipping). -
\textbf{Vận tốc góc (Gyro):} Gần 95\% thời gian các hoạt động, tốc độ
xoay người (RMS) chỉ nằm dưới 111 dps. Ngay cả ở mốc 99.9\% (những pha
lộn vòng hoặc vấp ngã mạnh nhất), tốc độ xoay cực đại cũng chỉ chạm
463.57 dps. Đây là \textbf{``bằng chứng thép''} cho thấy việc chọn thang
đo \texttt{±500\ dps} là chính xác tuyệt đối, tránh lãng phí độ phân
giải khi nén xuống INT8 (so với mốc ±2000 dps của các nghiên cứu cũ).

\textbf{Bảng Phân Tích Bách Phân Vị Theo Từng Hành Động (Đã Cắt
Window):}

\paragraph{1. Bảng Phân Tích Bách Phân Vị Gia Tốc (Accel) Theo Từng Hành
Động}\label{bux1ea3ng-phuxe2n-tuxedch-buxe1ch-phuxe2n-vux1ecb-gia-tux1ed1c-accel-theo-tux1eebng-huxe0nh-ux111ux1ed9ng}

\textbf{Phần 1: Trục X và Trục Y}

{\def\LTcaptype{none} % do not increment counter
\begingroup
\small
\setlength{\tabcolsep}{4pt}
\renewcommand{\arraystretch}{1.3}
\begin{longtable}{@{} l c c c c c c @{}}
\toprule\noalign{}
Nhãn (Class) & X-P50 & X-P95 & X-P99.9 & Y-P50 & Y-P95 & Y-P99.9 \\
\midrule\noalign{}
\endhead
\bottomrule\noalign{}
\endlastfoot
\textbf{Walk} & 0.10 & 0.39 & 0.80 & 0.95 & 1.48 & 2.18 \\
\textbf{Run} & 0.20 & 0.94 & 2.17 & 0.83 & 3.04 & 5.77 \\
\textbf{Idle\_Lie} & 0.82 & 1.00 & 1.18 & 0.20 & 0.70 & 0.97 \\
\textbf{Idle\_StandSit} & 0.04 & 0.22 & 0.84 & 0.97 & 1.01 & 1.37 \\
\textbf{Trans} & 0.07 & 0.95 & 1.29 & 0.85 & 1.05 & 3.27 \\
\textbf{Fall} & 0.25 & 1.31 & 5.20 & 0.45 & 1.36 & 5.79 \\
\end{longtable}
\endgroup
}

\textbf{Phần 2: Trục Z và Gia tốc Tổng hợp (SVM)}

{\def\LTcaptype{none} % do not increment counter
\begingroup
\small
\setlength{\tabcolsep}{4pt}
\renewcommand{\arraystretch}{1.3}
\begin{longtable}{@{} l c c c c c c @{}}
\toprule\noalign{}
Nhãn (Class) & Z-P50 & Z-P95 & Z-P99.9 & SVM-P50 & SVM-P95 & SVM-P99.9 \\
\midrule\noalign{}
\endhead
\bottomrule\noalign{}
\endlastfoot
\textbf{Walk} & 0.23 & 0.69 & 1.56 & 1.00 & 1.61 & 2.38 \\
\textbf{Run} & 0.27 & 1.36 & 4.64 & 1.03 & 3.34 & 6.07 \\
\textbf{Idle\_Lie} & 0.45 & 0.99 & 1.18 & 1.01 & 1.07 & 1.36 \\
\textbf{Idle\_StandSit} & 0.27 & 0.56 & 1.02 & 1.01 & 1.04 & 1.54 \\
\textbf{Trans} & 0.38 & 0.90 & 1.85 & 1.00 & 1.20 & 3.62 \\
\textbf{Fall} & 0.46 & 1.17 & 3.79 & 0.98 & 2.19 & 8.18 \\
\end{longtable}
\endgroup
}

\paragraph{2. Bảng Phân Tích Bách Phân Vị Vận Tốc Góc (Gyro) Theo Từng
Hành
Động}\label{bux1ea3ng-phuxe2n-tuxedch-buxe1ch-phuxe2n-vux1ecb-vux1eadn-tux1ed1c-guxf3c-gyro-theo-tux1eebng-huxe0nh-ux111ux1ed9ng}

\textbf{Phần 1: Trục X và Trục Y}

{\def\LTcaptype{none} % do not increment counter
\begingroup
\small
\setlength{\tabcolsep}{4pt}
\renewcommand{\arraystretch}{1.3}
\begin{longtable}{@{} l c c c c c c @{}}
\toprule\noalign{}
Nhãn (Class) & X-P50 & X-P95 & X-P99.9 & Y-P50 & Y-P95 & Y-P99.9 \\
\midrule\noalign{}
\endhead
\bottomrule\noalign{}
\endlastfoot
\textbf{Walk} & 11.29 & 53.71 & 159.85 & 15.20 & 65.55 & 145.94 \\
\textbf{Run} & 45.29 & 231.63 & 512.57 & 34.97 & 139.65 & 280.94 \\
\textbf{Idle\_Lie} & 2.01 & 27.16 & 87.45 & 4.33 & 28.56 & 74.34 \\
\textbf{Idle\_StandSit} & 1.40 & 30.33 & 109.62 & 4.09 & 21.00 &
95.40 \\
\textbf{Trans} & 13.37 & 67.87 & 263.82 & 6.47 & 46.08 & 108.01 \\
\textbf{Fall} & 22.03 & 143.07 & 401.67 & 25.63 & 177.73 & 403.14 \\
\end{longtable}
\endgroup
}

\textbf{Phần 2: Trục Z và Vận tốc góc Tổng hợp (RMS)}

{\def\LTcaptype{none} % do not increment counter
\begingroup
\small
\setlength{\tabcolsep}{4pt}
\renewcommand{\arraystretch}{1.3}
\begin{longtable}{@{} l c c c c c c @{}}
\toprule\noalign{}
Nhãn (Class) & Z-P50 & Z-P95 & Z-P99.9 & RMS-P50 & RMS-P95 & RMS-P99.9 \\
\midrule\noalign{}
\endhead
\bottomrule\noalign{}
\endlastfoot
\textbf{Walk} & 11.60 & 43.21 & 80.08 & 30.59 & 87.21 & 177.25 \\
\textbf{Run} & 29.42 & 90.76 & 203.80 & 83.98 & 267.63 & 532.29 \\
\textbf{Idle\_Lie} & 0.98 & 22.46 & 98.69 & 5.51 & 48.35 & 123.11 \\
\textbf{Idle\_StandSit} & 0.37 & 19.29 & 77.03 & 4.52 & 50.00 &
123.47 \\
\textbf{Trans} & 3.97 & 48.83 & 113.68 & 25.29 & 87.89 & 272.26 \\
\textbf{Fall} & 12.02 & 113.71 & 351.99 & 55.40 & 246.43 & 562.54 \\
\end{longtable}
\endgroup
}

\emph{Nhận xét bổ sung:} - \textbf{Về ngưỡng cắt 250 dps vs 500 dps:}
Bảng trên cho thấy rõ ở mốc cực đại (P99.9), tốc độ xoay của té ngã
(Fall) có thể lên tới \textbf{562 dps}, và đỉnh điểm của hành động chạy
(Run) là \textbf{532 dps}. Nếu chọn thang đo 250 dps, mô hình sẽ hoàn
toàn ``mù'' (bị bão hòa tín hiệu) trước mọi sự kiện có tốc độ cao, khiến
nó nhầm lẫn giữa một cú ngã mạnh và một bước chạy nhanh. - \textbf{Tách
biệt các nhóm hành động tĩnh:} Nhóm \texttt{Idle\_Lie} và
\texttt{Idle\_StandSit} có cường độ gia tốc và vận tốc góc cực kỳ thấp
(P95 chỉ quanh quẩn 1g và 50 dps). Việc duy trì độ phân giải ở 500 dps
(thay vì 2000) giúp mô hình TCN và ResNet bắt được những vi chấn động
nhỏ này mà không bị làm tròn về 0 trong chuẩn INT8.

\begin{center}\rule{0.5\linewidth}{0.5pt}\end{center}

\subsection{2. PHÂN TÍCH CHI TIẾT CÁC MÔ HÌNH ĐÃ HUẤN
LUYỆN}\label{phuxe2n-tuxedch-chi-tiux1ebft-cuxe1c-muxf4-huxecnh-ux111uxe3-huux1ea5n-luyux1ec7n}

\subsubsection{2.1. Nhóm mô hình CNN Thuần (Ưu tiên
ESP-NN)}\label{nhuxf3m-muxf4-huxecnh-cnn-thuux1ea7n-ux1b0u-tiuxean-esp-nn}

Đây là nhóm mô hình hoàn toàn tuân thủ các quy tắc phần cứng của vi điều
khiển ESP32-S3 (ESP-NN), chỉ sử dụng các toán tử được tăng tốc phần cứng
như \texttt{Conv2D}, \texttt{DepthwiseConv2D}, \texttt{MaxPool},
\texttt{Mean}, và \texttt{ReLU/ReLU6}.

{\def\LTcaptype{none} % do not increment counter
\begingroup
\small
\setlength{\tabcolsep}{4pt}
\renewcommand{\arraystretch}{1.3}
\begin{longtable}{@{} l p{5cm} c c c @{}}
\toprule\noalign{}
Tên mô hình & Cấu trúc \& Kỹ thuật & Kích thước INT8 & Tốc độ Firmware (SRAM) & Fall Recall (INT8) \\
\midrule\noalign{}
\endhead
\bottomrule\noalign{}
\endlastfoot
\textbf{v30} & CNN thuần 6 trục. Sử dụng Separable Conv1D. Lược bỏ hoàn
toàn Sigmoid và SE Block. & \textasciitilde25.0 KB & \textbf{19.78 ms} &
\textbf{98.00\%} \\
\textbf{v31} & CNN thuần 4 trục (Accel + Gyro RMS). Giống v30 nhưng đầu
vào nhỏ hơn. Nhẹ nhất. & \textasciitilde24.9 KB & \textbf{19.25 ms} &
\textbf{98.00\%} \\
\end{longtable}
\endgroup
}

\emph{Đánh giá:} CNN thuần là lựa chọn tối ưu tuyệt đối cho phần cứng.
Tốc độ thực thi chớp nhoáng (\textless{} 20ms) cho phép triển khai
Real-time với tần số quét cực cao mà không lo quá tải CPU.

\begin{center}\rule{0.5\linewidth}{0.5pt}\end{center}

\subsubsection{2.2. Nhóm mô hình Lai (CNN +
LSTM)}\label{nhuxf3m-muxf4-huxecnh-lai-cnn-lstm}

{\def\LTcaptype{none} % do not increment counter
\begingroup
\small
\setlength{\tabcolsep}{4pt}
\renewcommand{\arraystretch}{1.3}
\begin{longtable}{@{} l p{5cm} c c c @{}}
\toprule\noalign{}
Tên mô hình & Cấu trúc \& Kỹ thuật & Kích thước INT8 & Tốc độ Firmware (SRAM) & Fall Recall (INT8) \\
\midrule\noalign{}
\endhead
\bottomrule\noalign{}
\endlastfoot
\textbf{v30\_lstm32} & Trích xuất đặc trưng bằng CNN, sau đó đưa qua 1
lớp LSTM (32 units). Phải thiết lập Keras 2 Legacy để ép TFLite gom
thành op \texttt{UNIDIRECTIONAL\_SEQUENCE\_LSTM}. & \textasciitilde39.2
KB & \textbf{210.95 ms} & 94.50\% \\
\end{longtable}
\endgroup
}

\emph{Đánh giá:} Vì ESP-NN không hỗ trợ tăng tốc bằng vector phần cứng
cho LSTM, vi điều khiển phải chạy bằng code C++ thuần. Tốc độ giảm gấp
10 lần so với CNN.

\begin{center}\rule{0.5\linewidth}{0.5pt}\end{center}

\subsubsection{2.3. Nhóm mô hình
ResNet1D}\label{nhuxf3m-muxf4-huxecnh-resnet1d}

{\def\LTcaptype{none} % do not increment counter
\begingroup
\small
\setlength{\tabcolsep}{4pt}
\renewcommand{\arraystretch}{1.3}
\begin{longtable}{@{} l p{5cm} c c c @{}}
\toprule\noalign{}
Tên mô hình & Cấu trúc \& Kỹ thuật & Kích thước INT8 & Tốc độ Firmware (SRAM) & Fall Recall (INT8) \\
\midrule\noalign{}
\endhead
\bottomrule\noalign{}
\endlastfoot
\textbf{v30\_resnet1d} & ResNet1D tích hợp Squeeze-and-Excitation (SE)
Block. & \textasciitilde83.3 KB & \textbf{45.75 ms} & 96.00\% \\
\end{longtable}
\endgroup
}

\emph{Đánh giá:} Mặc dù là CNN, nhưng việc tồn tại hàm \texttt{Sigmoid}
(trong khối SE Block) và \texttt{GlobalMaxPooling}
(\texttt{REDUCE\_MAX}) khiến một phần mạng không được tăng tốc. Ngoài
ra, việc chứa nhiều tensor nhỏ rải rác khiến lượng Metadata (thông điệp
đi kèm) phình to (file tflite nặng tới 83KB dù số lượng tham số ít).

\textbf{Giải mã sự suy giảm tốc độ của v30 so với v25 (45ms vs 32ms):}
Dù \texttt{v30\_resnet1d} giữ nguyên kiến trúc gốc giống hệt
\texttt{v25} (cùng version Keras, cùng thuật toán nén INT8), tốc độ suy
luận thực tế trên ESP32 vẫn bị chậm đi (từ 32ms lên 45.75ms). Khi ``mổ
xẻ'' các tensor trong file \texttt{.tflite}, nguyên nhân cốt lõi đã lộ
diện ở tầng lượng tử hóa (Quantization): - Việc v30 phải học trên tập dữ
liệu \texttt{Trans} khó hơn (có augmentation) khiến mạng phải thay đổi
mạnh phân bố trọng số để hội tụ. Hệ quả là các dải Lượng tử hóa
(\texttt{scale} và \texttt{zero\_point}) của các lớp Conv ở v30 bị thay
đổi hoàn toàn so với v25 (Ví dụ: \texttt{out\_scale} của lớp Conv2D số 2
bị lệch từ \texttt{0.1511} ở v25 xuống \texttt{0.1257} ở v30). - Thư
viện tăng tốc ESP-NN cực kỳ nhạy cảm với dải \texttt{scale}. Nó cần tính
toán một biến
\texttt{quantized\_multiplier\ =\ (in\_scale\ *\ filter\_scale)\ /\ out\_scale}.
Nếu tỷ lệ này vượt ra khỏi giới hạn các lệnh dịch bit (bit-shift) tối ưu
của vi điều khiển, ESP-NN sẽ tự động ngắt tính năng Hardware
Acceleration và đẩy lớp Conv đó về chạy bằng vòng lặp C++ thuần
(Reference Kernel) nhằm đảm bảo không bị tràn số (overflow). Chính sự
thay đổi trọng số ngẫu nhiên do dữ liệu khó đã vô tình đẩy một vài lớp
Conv của v30 rơi vào trường hợp ``cấm'' này, tạo ra độ trễ
\textasciitilde13ms.

\begin{center}\rule{0.5\linewidth}{0.5pt}\end{center}

\subsubsection{2.4. Nhóm mô hình TCN (Temporal Convolutional
Network)}\label{nhuxf3m-muxf4-huxecnh-tcn-temporal-convolutional-network}

TCN sử dụng kiến trúc tích chập theo thời gian, thường dùng Dilated
Convolutions (Tích chập giãn nở) để tăng trường nhìn (Receptive Field)
mà không làm tăng tham số.

{\def\LTcaptype{none} % do not increment counter
\begingroup
\small
\setlength{\tabcolsep}{4pt}
\renewcommand{\arraystretch}{1.3}
\begin{longtable}{@{} l p{5cm} c c c @{}}
\toprule\noalign{}
Tên mô hình & Cấu trúc \& Kỹ thuật & Kích thước INT8 & Tốc độ Firmware (SRAM) & Fall Recall (INT8) \\
\midrule\noalign{}
\endhead
\bottomrule\noalign{}
\endlastfoot
\textbf{v30\_tcn} & TCN Cổ điển. Sử dụng Dilated Convolutions (Dilation
\textgreater{} 1). & \textasciitilde105.0 KB & \textbf{2058.50 ms} &
100.0\% (Mẫu nhỏ) \\
\textbf{v30\_tcn\_optimize} & TCN thiết kế lại: Bỏ Dilation, thay bằng
chuỗi \texttt{DepthwiseConv2D} và \texttt{Conv2D} chuẩn. &
\textasciitilde66.5 KB & \textbf{184.98 ms} & 96.50\% \\
\textbf{v30\_tcn\_optimize\_v2} & Tối ưu lại Stride và Kernel Size so
với v1 để tìm điểm cân bằng. & \textasciitilde65.6 KB & \textbf{217.75
ms} & 96.50\% \\
\end{longtable}
\endgroup
}

\emph{Đánh giá:} - \textbf{TCN Cổ điển} là ``thảm hoạ'' phần cứng: Mất
hơn 2 giây để dự đoán 1 mẫu. Lý do là ESP-NN từ chối tăng tốc bất kỳ
Convolution nào có \texttt{Dilation\_rate\ \textgreater{}\ 1}. -
\textbf{TCN Optimize} đã chứng minh việc thiết kế mạng thuận theo phần
cứng (Hardware-aware Neural Architecture Design) là bắt buộc. Việc đổi
sang Conv chuẩn giúp tốc độ tăng hơn 10 lần (184ms vs 2058ms).

\begin{center}\rule{0.5\linewidth}{0.5pt}\end{center}

\subsection{3. KẾT LUẬN \& ĐỀ XUẤT CHO TRIỂN KHAI THỰC
TẾ}\label{kux1ebft-luux1eadn-ux111ux1ec1-xuux1ea5t-cho-triux1ec3n-khai-thux1ef1c-tux1ebf}

\begin{enumerate}
\def\labelenumi{\arabic{enumi}.}
\tightlist
\item
  \textbf{Hiệu suất (Inference Speed):} ESP32-S3 chỉ toả sáng khi mạng
  \textbf{tuân thủ nghiêm ngặt chuẩn ESP-NN} (Chỉ dùng Conv chuẩn,
  Depthwise, MaxPool, Mean, ReLU). Bất kỳ toán tử ngoại lai nào (LSTM,
  Sigmoid, Dilated Conv, Reduce\_Max) đều sẽ kéo tốc độ tụt lùi từ 3 đến
  100 lần.
\item
  \textbf{Kích thước Flash/RAM:} Mô hình có thiết kế ít khối (Layer)
  nhưng kích thước lớn sẽ tiết kiệm bộ nhớ hơn là mô hình ít tham số
  nhưng xé lẻ thành hàng chục khối nhỏ (như SE Block của ResNet1D), do
  overhead từ FlatBuffer Metadata của TFLite.
\item
  \textbf{Đề xuất Deploy:} Phiên bản \textbf{\texttt{v31} (CNN 4 trục)}
  và \textbf{\texttt{v30} (CNN 6 trục)} là 2 ứng cử viên sáng giá nhất
  cho sản phẩm cuối cùng. Chúng đạt độ trễ tuyệt hảo
  (\textasciitilde19ms), dung lượng siêu gọn (\textasciitilde25KB) mà
  vẫn giữ vững tỷ lệ bắt té ngã cực kỳ ấn tượng (98\%). Sự đánh đổi
  (Trade-off) bằng các mô hình phức tạp hơn như TCN hay LSTM là không
  mang lại đủ lợi ích so với lượng độ trễ phát sinh.
\end{enumerate}

\end{document}
